\documentclass[11pt]{article}
\usepackage{fontspec}
\usepackage[french,english]{babel}
\usepackage{amsmath}
\usepackage{amssymb}
\usepackage{amsthm}
\usepackage{mathrsfs}
\usepackage{tikz-cd}
\usepackage{enumitem}

\newcommand{\printedpage}[1]{\marginpar{\footnotesize p.~#1}}
\newcommand{\msp}[1]{\marginpar{\footnotesize [#1]}}
\newcommand{\circled}[1]{\tikz[baseline=(c.base)]\node[draw,circle,inner sep=1pt](c){#1};}
\newcommand{\ovalbox}[1]{\tikz[baseline=(c.base)]\node[draw,rounded corners=6pt,inner sep=2pt](c){#1};}
\newcommand{\Ob}{\operatorname{Ob}}
\newcommand{\card}{\operatorname{card}}
\newsavebox{\dctauAB}
\newsavebox{\dctauABone}
\newsavebox{\dctauAAzero}

\begin{document}

\sbox{\dctauAB}{\circled{$\scriptstyle\tau_{\alpha,\beta}$}}% boxes for the circled labels of ms. 318 (p. 120)
\sbox{\dctauABone}{\circled{$\scriptstyle\tau_{\alpha,\beta+1}$}}
\sbox{\dctauAAzero}{\circled{$\scriptstyle\tau_{\alpha,\alpha_0}$}}
\printedpage{119}
% head: A. Grothendieck : Dérivateurs, Ch. XVIII, Édition des Universités de Montpellier II et Paris VII
satisfying the constraints $(**)$, $(***)$, $(****)$. We are going to make the construction in three steps.

\circled{1} Start:
\[
i_0 \;=\; i\ .
\]
Since $J'$ is cofinal in $J$, exists $j_0\in\Ob J'$ and
\[
\varphi(i_0)\xrightarrow{u_0} j_0\ .
\]

\msp{page 316}

\circled{2} Suppose that the $i_\beta$, $j_\beta$ have been constructed for $\beta<\alpha$, as well as the $\tau_{\beta',\beta}:i_\beta\longrightarrow i_{\beta'}$, $\tau'_{\beta',\beta}:j_\beta\longrightarrow j_{\beta'}$ for $\beta,\beta'<\alpha$, finally the $u_\beta:\varphi(i_\beta)\longrightarrow j_\beta$, $v_\beta:j_\beta\longrightarrow\varphi(i_{\beta+1})$ for $\beta<\alpha$ resp. for $\beta+1<\alpha$, so as to satisfy the constraints $(**)$, $(****)$, i.e. commutativities in
\[
\left\{
\begin{array}{l}
\left.
\begin{array}{ll}
\circled{1} &
\begin{tikzcd}[ampersand replacement=\&]
\varphi(i_\beta) \arrow[r, "u_\beta"] \arrow[rr, bend right=15, "\varphi(\tau_{\beta+1,\beta})"'] \& j_\beta \arrow[r, "v_\beta"] \& \varphi(i_{\beta+1})
\end{tikzcd}
\\[5ex]
\circled{2} &
\begin{tikzcd}[ampersand replacement=\&]
j_\beta \arrow[r, "v_\beta"] \arrow[rr, dotted, bend right=15, "\tau'_{\beta+1,\beta}"'] \& \varphi(i_{\beta+1}) \arrow[r, "u_{\beta+1}"] \& j_{\beta+1}
\end{tikzcd}
\end{array}
\right\}\quad \textit{if}\ \beta+1<\alpha
\\[8ex]
\begin{array}{ll}
\circled{3} &
\begin{tikzcd}[ampersand replacement=\&]
\varphi(i_\beta) \arrow[r, "u_\beta"] \arrow[d, "\varphi(\tau_{\beta',\beta})"'] \& j_\beta \arrow[d, "\tau'_{\beta',\beta}"] \\
\varphi(i_{\beta'}) \arrow[r, "u_{\beta'}"] \& j_{\beta'}
\end{tikzcd}
\qquad\qquad \text{if } \beta<\beta'<\alpha
\end{array}
\\[10ex]
\begin{array}{ll}
\circled{4} &
\begin{tikzcd}[ampersand replacement=\&]
j_\beta \arrow[r, "v_\beta"] \arrow[d, "\tau'_{\beta',\beta}"'] \& \varphi(i_{\beta+1}) \arrow[d, "\varphi(\tau_{\beta'+1,\beta+1})"] \\
j_{\beta'} \arrow[r, "v_{\beta'}"] \& \varphi(i_{\beta'+1})
\end{tikzcd}
\qquad\qquad \text{if } \beta<\beta' \text{ and } \beta'+1<\alpha\ ,
\end{array}
\end{array}
\right.
\]

\msp{page 317}

we must construct
% foot: Version 17 décembre 2025

\printedpage{120}
% head: A. Grothendieck : Dérivateurs, Ch. XVIII, Édition des Universités de Montpellier II et Paris VII
\begin{itemize}[leftmargin=*]
\item[1)] $i_\alpha\in\Ob I$,\quad and a homomorphism $(i_\beta)_{\beta\in I_\alpha}\longrightarrow i_\alpha$, i.e. $\tau_{\alpha,\beta}$ : $i_\beta\longrightarrow i_\alpha$ $(\beta\in I_\alpha)$, compatible with the translations $\tau_{\beta',\beta}$ for $\beta<\beta'<\alpha$.
\item[2)] $j_\alpha\in\Ob J'$,\quad and a homomorphism $(j_\beta)_{\beta\in I_\alpha}\longrightarrow j_\alpha$, i.e. $\tau'_{\alpha,\beta}$ : $j_\beta\longrightarrow j_\alpha$ $(\beta\in I_\alpha)$, compatible with the translations $\tau'_{\beta',\beta}$ for $\beta<\beta'<\alpha$.
\item[3)] An arrow
\[
u_\alpha:\varphi(i_\alpha)\longrightarrow j_\alpha\ .
\]
\item[4)] If $\alpha$ is not a limit ordinal, i.e. $\alpha=\alpha_0+1$, an arrow
\[
v_{\alpha_0}:j_{\alpha_0}\longrightarrow\varphi(i_{\alpha_0+1})=\varphi(i_\alpha)\ ;
\]
\end{itemize}

and this so as to satisfy the constraints \circled{1}\,\circled{2}\,\circled{3}\,\circled{4} when we assume $\beta+1\leqslant\alpha$ (instead of $\beta'<\alpha$) in \circled{1} and \circled{2}, and $\beta'\leqslant\alpha$ (instead of $\beta'<\alpha$) in \circled{3}, and $\beta'+1\leqslant\alpha$ (instead of $\beta'+1<\alpha$) in \circled{4}. So for the constraints $1^\circ$, $2^\circ$, $4^\circ$, we need take them into account \textit{only} if $\alpha$ is not a limit ordinal, so $\alpha=\alpha_0+1$, and setting there % [sic] il n'y a en tenir -> il n'y a à en tenir, the preposition « à » appears to be missing
\[
\beta' \;=\; \alpha_0\ .
\]
As for the constraint $3^\circ$, we must take it
\msp{page 318}
into account for $\beta'=\alpha$. So in all, we must require
\[
(3_\alpha)(\beta)\qquad\qquad
\begin{tikzcd}
\varphi(i_\beta) \arrow[r, "u_\beta"] \arrow[d, "\varphi(\usebox{\dctauAB})"'] & j_\beta \arrow[d, "\tau'_{\alpha,\beta}" {draw, circle, inner sep=1pt}] \\
\varphi(i_\alpha) \arrow[r, "u_\alpha" {draw, circle, inner sep=1pt}] & j_\beta % [sic] j_\beta -> j_\alpha, the target of \tau'_{\alpha,\beta} and u_\alpha is j_\alpha
\end{tikzcd}
\qquad \text{commutative for } \beta<\alpha\ (\footnotemark[140]),
\]
\footnotetext[140]{\textbf{N.B.} It suffices to require it for $\beta=\alpha_0$, in the case $\alpha=\alpha_0+1$, cf. below, p. 323.}
(by restriction on $\alpha$), and if $\alpha=\alpha_0+1$ is a non-limit ordinal, we must require in addition
\[
(4_\alpha)(\beta)\qquad\qquad
\begin{tikzcd}
j_\beta \arrow[r, "v_\beta"] \arrow[d, "\tau'_{\alpha_0,\beta}"'] & \varphi(i_{\beta+1}) \arrow[d, "\varphi(\usebox{\dctauABone})"] \\
j_{\alpha_0} \arrow[r, "v_{\alpha_0}" {draw, circle, inner sep=1pt}] & \varphi(i_{\alpha_0+1})=\varphi(i_\alpha)
\end{tikzcd}
\qquad \beta<\alpha_0
\]
\[
(1_\alpha)\qquad\qquad
\begin{tikzcd}
\varphi(i_{\alpha_0}) \arrow[r, "u_{\alpha_0}"] \arrow[rr, bend right=15, "\varphi(\usebox{\dctauAAzero})"'] & j_{\alpha_0} \arrow[r, "v_{\alpha_0}" {draw, circle, inner sep=1pt}] & \varphi(i_{\alpha_0+1})=\varphi(i_\alpha)
\end{tikzcd}
\]
% foot: Version 17 décembre 2025

\printedpage{121}
% head: A. Grothendieck : Dérivateurs, Ch. XVIII, Édition des Universités de Montpellier II et Paris VII
\[
(2_\alpha)\qquad\qquad
\begin{tikzcd}
j_{\alpha_0} \arrow[r, "v_{\alpha_0}" {draw, circle, inner sep=1pt}] \arrow[rr, bend right=15, "\tau'_{\alpha,\alpha_0}"' {draw, circle, inner sep=1pt}] & \varphi(i_{\alpha_0+1})=\varphi(i_\alpha) \arrow[r, "u_\alpha" {draw, circle, inner sep=1pt}] & j_\alpha\ .
\end{tikzcd}
\]

In these compatibilities, I have circled the arrows $\tau_{\alpha,\beta}$, $\tau'_{\alpha,\beta}$ $(\beta<\alpha)$, $\Bigl(\ \tau_{\alpha,\beta+1}\ (\beta+1<\alpha)\ \Bigr)$, $v_{\alpha_0}$, $\Bigl(\ \tau_{\alpha,\alpha_0}\ \Bigr)$, $u_\alpha$, $\Bigl(\ \tau'_{\alpha,\alpha_0}\ \Bigr)$, which were introduced in the data 1) 2) 3) 4), page 317. % [sic] donnés -> données, the noun « données » is meant

\msp{page 319}

We are going to distinguish then the two cases.

\textit{2 A)} $\alpha=\alpha_0+1$.

To construct $i_\alpha$ and the $\tau_{\alpha,\beta}:i_\beta\longrightarrow i_\alpha$, it suffices to construct $i_\alpha$ and
\[
\circled{$\tau_{\alpha,\alpha_0}$}\ :\ i_{\alpha_0}\ \longrightarrow\ i_\alpha=i_{\alpha_0+1}\ .
\]
Similarly for $j_\alpha$ and the $\tau'_{\alpha,\beta}$, it suffices to construct
\[
\circled{$\tau'_{\alpha,\alpha_0}$}\ :\ j_{\alpha_0}\ \longrightarrow\ j_\alpha\ .
\]
We must moreover give ourselves
\[
j_{\alpha_0}\ \xrightarrow{\circled{$v_{\alpha_0}$}}\ \varphi(\underbrace{\,i_\alpha\,}_{=i_{\alpha_0+1}})\ \xrightarrow{\circled{$u_\alpha$}} j_\alpha\ ,
\]
so in all two objects $i_\alpha$, $j_\alpha$ (in $I$ and in $J'$ respectively) and four arrows $\tau_{\alpha,\alpha_0}$, $\tau'_{\alpha,\alpha_0}$, $v_{\alpha_0}$, $u_\alpha$. We need four compatibilities (page 318), where the last two $3_\alpha(\beta)$, $4_\alpha(\beta)$ depend on a parameter $\beta<\alpha$, i.e. $\beta\leqslant\alpha_0$, or on $\beta<\alpha_0$.

\msp{page 320 struck out}

\msp{page 321}

We are first going to deal with the first two compatibilities $1_\alpha$, $2_\alpha$, in terms of the datum of
\[
i_{\alpha_0}\ \in\ \Ob I\ ,\qquad j_{\alpha_0}\ \in\ \Ob J'\ ,\qquad u_{\alpha_0}\ :\ \varphi(i_{\alpha_0})\ \longrightarrow\ j_{\alpha_0}\ .
\]
The two desired compatibilities for the two objects $i_\alpha$, $j_\alpha$ and the four arrows $\tau_{\alpha,\alpha_0}$, $\tau'_{\alpha,\alpha_0}$, $v_{\alpha_0}$, $u_\alpha$ are indicated in the diagram
\[
(*)\qquad\qquad
\begin{tikzcd}
\varphi(i_{\alpha_0}) \arrow[r, "u_{\alpha_0}"] \arrow[rr, bend right=15, "\varphi(\tau_{\alpha,\alpha_0})"'] & j_{\alpha_0} \arrow[r, "v_{\alpha_0}\ ?"] \arrow[rr, bend left=20, "\tau'_{\alpha,\alpha_0}"] & \varphi(\underbrace{\,i_\alpha\,}_{?}) \arrow[r, "u_\alpha\ ?"] & j_\alpha\ .
\end{tikzcd}
\]

We note that since $\varphi$ is cofinal, $\exists\ i_\alpha=i\in\Ob I$ and $j_{\alpha_0}\xrightarrow{v_{\alpha_0}}\varphi(i)$. But we want the composite
\[
w\ :\ \varphi(i_{\alpha_0})\ \xrightarrow{u_{\alpha_0}}\ j_{\alpha_0}\ \xrightarrow{v_{\alpha_0}}\ \varphi(i)
\]
to be of the form $\varphi(\tau)$, where $\tau=\tau_{\alpha,\alpha_0}:i_{\alpha_0}\longrightarrow i_\alpha$ is an arrow in $I$. Since $\varphi$ is a cofinal functor of filtering categories, we know
% foot: Version 17 décembre 2025
\printedpage{122}
% head: A. Grothendieck : Dérivateurs, Ch. XVIII, Édition des Universités de Montpellier II et Paris VII
\msp{page 322}
that we can find $i\xrightarrow{\rho}i'$ in $I$, such that the composite arrow $\varphi(i_{\alpha_0})\xrightarrow{w}\varphi(i)\xrightarrow{\varphi(\rho)}\varphi(i')$ is of the form $\varphi(\tau)$. So at the cost of replacing $i$ by $i'$, and $v_{\alpha_0}$ by the composite $v'_{\alpha_0}$ : $j_{\alpha_0}\xrightarrow{v_{\alpha_0}}\varphi(i)\xrightarrow{\varphi(\rho)}\varphi(i')$, we can suppose that $v_{\alpha_0}u_{\alpha_0}$ is of the form $\varphi(\tau)$. This thus gives us in principle the object $i=i_\alpha$, and the arrows $v_{\alpha_0}:j_{\alpha_0}\longrightarrow\varphi(i)$ and $\tau_{\alpha,\alpha_0}:i_{\alpha_0}\longrightarrow i_\alpha$, so as to satisfy the first compatibility in $(*)$. It remains to construct $j_\alpha$, $\tau'_{\alpha,\alpha_0}$, $u_\alpha$ so as to satisfy the second compatibility $(*)$. Now $J'$ being cofinal in $J$, there exists $j_\alpha$ in $J'$ and an arrow $u_\alpha:\varphi(i_\alpha)\longrightarrow j_\alpha$. On the other hand, since $J'\longrightarrow J$ is fully faithful, the composite arrow $j_{\alpha_0}\xrightarrow{v_{\alpha_0}}\varphi(i_\alpha)\xrightarrow{u_\alpha}j_\alpha$ in $J$, between objects of $J'$, comes from a unique arrow $\tau'_{\alpha,\alpha_0}:j_{\alpha_0}\longrightarrow j_\alpha$ in $J'$.

\msp{page 323}

Thus, we have constructed the two objects and the four desired arrows, so as to satisfy the constraints $1_\alpha$, $2_\alpha$ of page 318. Are the constraints $3_\alpha(\beta)$, $4_\alpha(\beta)$ then satisfied automatically?
\[
\begin{tikzcd}
\varphi(i_\beta) \arrow[r, "u_\beta"] \arrow[d] & j_\beta \arrow[d] \\
\varphi(i_{\alpha_0}) \arrow[r, "u_{\alpha_0}"] \arrow[d, "\varphi(\tau_{\alpha,\alpha_0})"'] & j_{\alpha_0} \arrow[d, "\tau'_{\alpha,\alpha_0}"] \\
\varphi(i_\alpha) \arrow[r, "u_\alpha"] & j_\alpha
\end{tikzcd}
\qquad\qquad \beta<\alpha, \text{ i.e. } \beta\leqslant\alpha_0\ .
\]
The upper square is commutative by the induction hypothesis, for the total (composite) square to be commutative, it suffices that the lower square be so. We thus have a double arrow $(\tau'_{\alpha,\alpha_0}u_{\alpha_0},u_\alpha\varphi(\tau_{\alpha,\alpha_0})):\varphi(i_\alpha)\rightrightarrows j_\alpha$, it may not be trivial, but we know, $J'$ being filtering, that we can equalize it by an arrow $j_\alpha\longrightarrow\bar{j}_\alpha$ in $J'$. So at the cost of replacing $j_\alpha$ by % [sic] \varphi(i_\alpha) -> \varphi(i_{\alpha_0}), both composites of the double arrow start at \varphi(i_{\alpha_0})
\msp{page 324}
$\bar{j}_\alpha$, $u_\alpha$ by the composite $\bar{u}_\alpha:\varphi(i_0)\longrightarrow j_\alpha\longrightarrow\bar{j}_\alpha$ and $\tau'_{\alpha,\alpha_0}$ by the composite $\bar{\tau}'_{\alpha,\alpha_0}:j_{\alpha_0}\longrightarrow j_\alpha\longrightarrow\bar{j}_\alpha$ (without touching $i_\alpha$, $v_{\alpha_0}$, $\tau_{\alpha_0,\alpha}$), we can obtain commutativity, whence $3_\alpha(\beta)$. % [sic] \varphi(i_0) -> \varphi(i_\alpha), u_\alpha has source \varphi(i_\alpha); [sic] \tau_{\alpha_0,\alpha} -> \tau_{\alpha,\alpha_0}, the arrow is written \tau_{\alpha,\alpha_0} everywhere else

There remains the constraint $4_\alpha(\beta)$. For this, let us consider the composite diagram
\[
\begin{tikzcd}
j_\beta \arrow[r, "v_\beta"] \arrow[d, "\tau'_{\alpha_0,\beta}"'] & \varphi(i_{\beta+1}) \arrow[d, "\varphi(\tau_{\alpha,\beta+1})"] \\
j_{\alpha_0} \arrow[r, "v_{\alpha_0}"] & \varphi(i_{\alpha_0+1})=\varphi(i_\alpha)\ .
\end{tikzcd}
\]

If $\alpha_0$ is not a limit ordinal, but of the form $\alpha_1+1$, we can again reduce to the case where $\beta=\alpha_1$, proceeding as above. But we can proceed in all cases, without such an artifice (\,!), by considering the double arrows
\[
(*)\qquad\qquad
\begin{tikzcd}
j_\beta \arrow[r, shift left, "\varphi(\tau_{\alpha,\beta+1})v_\beta"] \arrow[r, shift right, "v_{\alpha_0}\tau'_{\alpha_0,\beta}"'] & \varphi(i_\alpha)
\end{tikzcd}
\]
% foot: Version 17 décembre 2025
\printedpage{123}
% head: A. Grothendieck : Dérivateurs, Ch. XVIII, Édition des Universités de Montpellier II et Paris VII
deduced from the preceding square. Since
\msp{page 325}
the set of these double arrows is of cardinality $\card\alpha<\pi$, and $J$ is large relative to $\pi$ (it would suffice that it be so for the cardinals $\pi'<\pi$), we can find an arrow $\varphi(i_\alpha)\xrightarrow{u}j$ in $J$ which equalizes all these double arrows. Since $\varphi$ is cofinal, we can suppose that $j$ is of the form $\varphi(i)$. Since $\varphi$ is a cofinal functor of filtering categories, there then exists $i\xrightarrow{\rho}i'$ in $I$, such that the composite $\varphi(i_\alpha)\xrightarrow{u}\varphi(i)\xrightarrow{\varphi(\rho)}\varphi(i')$ is of the form $\varphi(\sigma)$, for $\sigma:i_\alpha\longrightarrow i'$. In all, we have found $\sigma:i_\alpha\longrightarrow i'$, such that $\varphi(\sigma)$ $(=\varphi(\rho)u):\varphi(i_\alpha)\longrightarrow\varphi(i')$ equalizes all the double arrows considered in $(*)$, p. 324. So it is necessary to replace $i_\alpha$ by $\bar{i}_\alpha=i'$, $v_{\alpha_0}$ by the composite $j_{\alpha_0}\longrightarrow\varphi(i_\alpha)\xrightarrow{\varphi(\sigma)}\varphi(\bar{i}_\alpha)$, $\tau_{\alpha,\alpha_0}$
\msp{page 326}
by the composite $\bar{\tau}_{\alpha,\alpha_0}:i_{\alpha_0}\xrightarrow{\tau_{\alpha,\alpha_0}}i_\alpha\xrightarrow{\sigma}\bar{i}_\alpha$.

To sum up, here is the recipe for constructing, first \ovalbox{$i_\alpha$, $v_{\alpha_0}$, $\tau_{\alpha,\alpha_0}$,}and \textit{then} \ovalbox{$j_\alpha$, $u_\alpha$, $\tau'_{\alpha,\alpha_0}$}.

We proceed first as on page 321, 322 to construct the first triple, so as to satisfy the first compatibility $(*)$ (p. 321), i.e. $1_\alpha$. We then ``correct'' this choice of $i_\alpha$, $v_{\alpha_0}$, $\tau_{\alpha,\alpha_0}$, by replacing $i_\alpha$ by $\bar{i}_\alpha$ for $\sigma:i_\alpha\longrightarrow\bar{i}_\alpha$, $v_{\alpha_0}$ by $\bar{v}_{\alpha_0}=\varphi(\sigma)v_{\alpha_0}$, and $\tau_{\alpha,\alpha_0}$ by $\bar{\tau}_{\alpha,\alpha_0}=\sigma\tau_{\alpha,\alpha_0}$, so as to obtain the compatibilities $4_\alpha(\beta)$ for $\beta<\alpha_0$, as we have just indicated. We will no longer touch these choices of $i_\alpha$, $v_{\alpha_0}$, $\tau_{\alpha,\alpha_0}$, in the rest of the construction. We construct $(j_\alpha,u_\alpha,\tau'_{\alpha,\alpha_0})$ so as to satisfy the second compatibility $(*)$ (p. 321), as
\msp{page 327}
indicated on page 322. Then, we correct this choice, by replacing $j_\alpha$ by $\bar{j}_\alpha$, with $\rho:j_\alpha\longrightarrow\bar{j}_\alpha$, and $u_\alpha$ by $\bar{u}_\alpha=\rho u_\alpha$, $\tau'_{\alpha,\alpha_0}$ by $\bar{\tau}'_{\alpha,\alpha_0}=\rho\tau'_{\alpha,\alpha_0}$, so as to satisfy the constraints $3_\alpha(\beta)$ of page 318.

This completes the construction in the case where $\alpha=\alpha_0+1$.

\textit{Case 2 B: $\alpha$ is a limit ordinal.}

So we must construct the data 1) 2) 3) of page 317 (to the exclusion of datum 4), so as to satisfy the sole constraint $3_\alpha(\beta)$.

We begin to construct $i_\alpha$ and the $\tau_{\beta,\alpha}$ in $I$. This is possible, for $I_\alpha=\{\beta\mid\beta<\alpha\}$ is of cardinality $\pi_\alpha<\pi$ – it therefore suffices to suppose $I$ large relative to the cardinals $\pi'<\pi$.

We then likewise construct $J_\alpha$ in $J'$ and the $\tau'_{\beta,\alpha}$ $(\beta<\alpha)$, using the % [sic] J_\alpha -> j_\alpha, the object of J' is meant
\msp{page 328}
fact that $J'$ is large relative to $\card\alpha=\card I_\alpha$. We must moreover construct an arrow
\[
u_\alpha\ :\ \varphi(i_\alpha)\ \longrightarrow\ j_\alpha\ .
\]
This is possible, for $J'$ being cofinal in $J$, there exists $j'$ in $J'$ and $\varphi(i_\alpha)\longrightarrow J'$, and $J'$ being filtering, there exists $\bar{j}_\alpha$ and arrows % [sic] \varphi(i_\alpha)\longrightarrow J' -> \varphi(i_\alpha)\longrightarrow j', the arrow goes to the object j'
$\begin{tikzcd}[ampersand replacement=\&, baseline=(\tikzcdmatrixname-1-1.base)]
j_\alpha \arrow[rr] \& \& j_\alpha \\
\& j' \arrow[ur] \&
\end{tikzcd}$, so at the cost of replacing $j_\alpha$ by $\bar{j}_\alpha$, the $\tau'_{\alpha,\beta}$ by $\bar{\tau}'_{\alpha,\beta}$, we do indeed find an arrow $u_\alpha:\varphi(i_\alpha)\longrightarrow j_\alpha$. We must then deal with the % [sic] target j_\alpha in the small diagram -> \bar{j}_\alpha, the two arrows go to the object \bar{j}_\alpha just introduced
% foot: Version 17 décembre 2025

\end{document}
